% Prospectos Solares. Presentacion del procedimiento.
%
% Version SB. Rama propia del guion, independiente de presentacion/.
%
% Compilar desde presentacion_sb/:
%   python figuras.py
%   pdflatex -output-directory=../outputs/presentacion_sb/build prospectos_solares_sb.tex
%   pdflatex -output-directory=../outputs/presentacion_sb/build prospectos_solares_sb.tex
%   python verificar.py
%   python powerpoint.py          (exporta entregables/...pptx)
%   copy ../outputs/presentacion_sb/build/prospectos_solares_sb.pdf ^
%        ../entregables/presentacion_prospectos_solares.pdf
%
% El PDF que se revisa es SIEMPRE el de entregables/. La carpeta build/ es el
% taller de pdflatex y no se entrega: dejar dos mazos completos rodando fue lo
% que hizo que se revisara una version que no era.
%
% El guion sigue el procedimiento tal como lo describe Samuel Blanco: donde se ha
% construido, el mapa de grillas, el modelo de similitud, la caracterizacion y
% seleccion por criterios del cliente, y la bajada al lote. Cada tramo entre dos
% pasos tiene lamina propia y nombra el resultado concreto que viaja de uno al
% otro, porque eso es lo que hace auditable la cadena.
%
% Regla de redaccion de esta presentacion: el lector no ve nunca un nombre de
% fichero, ni de columna, ni de modulo. La trazabilidad se sostiene nombrando la
% ETAPA que produce el dato y la ENTIDAD que lo publica, que es lo que un tercero
% puede verificar. El fichero donde el proyecto lo guarda no le sirve a nadie.
%
% Segunda regla: una idea por lamina. Lo demas lo lleva la figura. Una figura
% encogida no se lee proyectada, de modo que cuando una lamina no cabe se recorta
% el texto antes que reducir la figura.
%
% Toda cifra de este fichero esta verificada contra la salida real.

\documentclass[aspectratio=169, 11pt]{beamer}

\usepackage[utf8]{inputenc}
\usepackage[T1]{fontenc}
\usepackage[spanish, es-nodecimaldot]{babel}
\usepackage{graphicx}
\usepackage{booktabs}
\usepackage{array}
\usepackage{tikz}
\usetikzlibrary{arrows.meta, calc}
\usepackage{xcolor}
\usepackage{helvet}
\renewcommand{\familydefault}{\sfdefault}
\renewcommand{\arraystretch}{1.25}

\graphicspath{%
  {../outputs/presentacion_sb/figuras/}%
  {../outputs/presentacion_sb/capturas/}}
% ============================================================================
% Prospectos Solares. Sistema visual de la presentacion.
% ----------------------------------------------------------------------------
% La paleta, el isotipo y las proporciones se toman del reporte de lotes
% (reporte_predios/plantilla.py, modo claro), para que la presentacion y el
% entregable se lean como un mismo producto y no como dos piezas distintas.
%
% Reglas del sistema:
%   Color      un solo verde azulado de marca. El acento oscuro (acento) es el
%              que lleva texto y linea; el claro (marca) solo identifica.
%              bien, aviso y mal se reservan para estado, nunca para decorar.
%   Rejilla    margen lateral unico (\margenlat) y altura util unica
%              (\altoutil). Ninguna lamina cambia de proporcion.
%   Filete     un solo divisor de marca: 1,8 cm por 1,4 pt en acento. Las
%              separaciones internas usan la hairline de 0,6 pt en linea.
%   Tipografia una sola familia de palo seco, tres pesos de gris (tinta,
%              tinta2, tinta3) y jerarquia por tamano, no por color.
%   Cajas      toda caja de diagrama es un nodo de tikz con text width y
%              minimum height: la caja se dimensiona a partir del texto, de
%              modo que el texto nunca puede desbordarla.
%   Lectura    sin justificar y sin partir palabras.
% ============================================================================

\usetheme{default}
\usefonttheme{professionalfonts}
\setbeamertemplate{navigation symbols}{}

% --- Paleta ------------------------------------------------------------------
\definecolor{marca}{HTML}{18919C}        % identidad
\definecolor{acento}{HTML}{0F6A72}       % texto y linea de acento
\definecolor{acento2}{HTML}{0B5057}      % acento fuerte, enfasis
\definecolor{acentosuave}{HTML}{E3EFF0}  % fondo de acento
\definecolor{tinta}{HTML}{1E2829}        % texto principal
\definecolor{tinta2}{HTML}{49585A}       % texto secundario
\definecolor{tinta3}{HTML}{5C6B6D}       % texto terciario, notas
\definecolor{linea}{HTML}{DCE2E2}        % hairline
\definecolor{fondo}{HTML}{F4F6F6}        % fondo de pagina
\definecolor{superficie}{HTML}{EAEFF0}   % superficie elevada
\definecolor{bien}{HTML}{2F6B4C}         % estado favorable
\definecolor{aviso}{HTML}{8A5C17}        % estado en curso
\definecolor{mal}{HTML}{9C3A1F}          % estado que impide

% Clasificacion, misma escala que el reporte de lotes.
\definecolor{clasei}{HTML}{17563E}
\definecolor{claseii}{HTML}{2C5F9E}
\definecolor{claseiii}{HTML}{C08322}
\definecolor{claseiv}{HTML}{879596}

\setbeamercolor{normal text}{fg=tinta, bg=white}
\setbeamercolor{background canvas}{bg=white}
\setbeamercolor{frametitle}{fg=tinta, bg=white}
\setbeamercolor{framesubtitle}{fg=tinta2}
\setbeamercolor{structure}{fg=acento}
\setbeamercolor{alerted text}{fg=acento2}

\setbeamerfont{frametitle}{size=\large, series=\bfseries}
\setbeamerfont{framesubtitle}{size=\scriptsize, series=\mdseries}

% --- Rejilla fija ------------------------------------------------------------
% Todas las laminas comparten el mismo margen y la misma altura util, para que
% las proporciones no bailen entre una y otra.
\newlength{\margenlat}\setlength{\margenlat}{1.0cm}
\newlength{\altoutil}\setlength{\altoutil}{0.615\paperheight}

% Sin partir palabras: en laminas con columnas estrechas el corte silabico
% ensucia la lectura y en una presentacion no compensa.
\hyphenpenalty=10000
\exhyphenpenalty=10000
\setbeamersize{text margin left=\margenlat, text margin right=\margenlat}

% --- Grosores del sistema ----------------------------------------------------
\newlength{\anchofilete}\setlength{\anchofilete}{1.8cm}
\newlength{\altofilete}\setlength{\altofilete}{1.4pt}
\newlength{\hairline}\setlength{\hairline}{0.6pt}

% Divisor de marca. Es la unica linea gruesa de toda la presentacion.
\newcommand{\reglamarca}[1][acento]{\textcolor{#1}{\rule{\anchofilete}{\altofilete}}}

% Hairline de ancho libre, para separar bloques dentro de una lamina.
\newcommand{\filete}[1][\linewidth]{\textcolor{linea}{\rule{#1}{\hairline}}}

% Bandera moderada. \raggedright deja estiramiento infinito al margen derecho y
% en columna estrecha produce lineas de longitudes muy dispares. Acotando el
% estiramiento a dos cuadratines las lineas se igualan bastante sin justificar,
% que es lo que el sistema no quiere porque abre rios y partiria palabras.
\newcommand{\bandera}{\rightskip=0pt plus 2em\relax\parindent=0pt}

% Rotulo breve sobre un bloque. Siempre en versales y en gris terciario.
\newcommand{\rotulo}[2][tinta3]{{\scriptsize\textcolor{#1}{\MakeUppercase{#2}}}}

% --- Isotipo: siete barras de altura creciente, extremos redondeados ---------
\newcommand{\isotipo}[2][marca]{%
  \begin{tikzpicture}[y=#2, x=#2]
    \foreach \x/\y/\h in {0/0.270/0.460, 0.152/0.215/0.570, 0.304/0.160/0.680,
                          0.457/0.110/0.780, 0.609/0.060/0.880, 0.761/0.025/0.950,
                          0.909/0/1.0}
      {\fill[#1, rounded corners=0.0456*#2] (\x, \y) rectangle (\x+0.0913, \y+\h);}
  \end{tikzpicture}%
}

\newcommand{\marcaTexto}[1][marca]{%
  \raisebox{-0.28em}{\isotipo[#1]{0.62em}}\hspace{0.5em}%
  \textcolor{#1}{\textbf{\footnotesize MÉTODOS MIXTOS}}\,%
  \textcolor{tinta3}{\scriptsize CONSULTORES}%
}

% --- Titulo de lamina --------------------------------------------------------
% Titulo, subtitulo opcional y el filete de marca. Las medidas verticales son
% las mismas en todas las laminas: el cuerpo siempre arranca a la misma altura.
\makeatletter
\setbeamertemplate{frametitle}{%
  \vspace{0.45em}%
  \begin{beamercolorbox}[wd=\paperwidth, leftskip=\margenlat, rightskip=\margenlat]{frametitle}%
    % Sin \raggedright: su definicion pone \leftskip a cero y con el se llevaba
    % la sangria del propio beamercolorbox. Los titulos salian pegados al borde
    % del papel mientras el cuerpo arrancaba en el margen. Se compone en bandera
    % fijando las dos sangrias a mano.
    % El \rightskip se fija con \dimexpr: \margenlat es un registro de glue y
    % detras de un registro TeX no lee el "plus", que se imprimia como texto.
    \leftskip=\margenlat
    \rightskip=\dimexpr\margenlat\relax plus 1fil\relax
    \parindent=0pt%
    \usebeamerfont{frametitle}\insertframetitle%
    \ifx\insertframesubtitle\@empty\else%
      \\[0.25em]{\usebeamerfont{framesubtitle}\usebeamercolor[fg]{framesubtitle}\insertframesubtitle}%
    \fi%
  \end{beamercolorbox}%
  \vspace{0.5em}%
  % Sin \hspace: fuera del beamercolorbox ya rige el margen de lamina, de modo
  % que la sangria se sumaba dos veces y el filete arrancaba a 2 cm mientras el
  % titulo y el cuerpo arrancaban a 1 cm.
  \reglamarca%
  % Aire bajo el filete. Estaba en 0,85 em y dejaba una banda muerta entre el
  % filete y la figura justo donde la figura tenia que empezar, lo que obligaba a
  % encoger despues. Con 0,35 em la lamina recupera alto util y varias figuras
  % bajan de shrink.
  \vspace{0.35em}%
}
\makeatother

% --- Pie ---------------------------------------------------------------------
% Seccion a la izquierda, isotipo y numero a la derecha. Sin linea: el pie se
% separa por aire, no por trazo.
\setbeamertemplate{footline}{%
  \hbox{%
  \begin{beamercolorbox}[wd=0.5\paperwidth, ht=2.2ex, dp=1.6ex, leftskip=\margenlat]{}%
    \tiny\textcolor{tinta3}{\MakeUppercase{\insertsectionhead}}%
  \end{beamercolorbox}%
  \begin{beamercolorbox}[wd=0.5\paperwidth, ht=2.2ex, dp=1.6ex, rightskip=\margenlat]{}%
    \hfill\raisebox{-0.15em}{\isotipo[linea]{0.44em}}\hspace{0.5em}\tiny\textcolor{tinta3}{\insertframenumber}%
  \end{beamercolorbox}}%
}

% Sin justificar: en columnas estrechas la justificacion abre rios de espacio y
% parte palabras. La alineacion se marca dentro de cada minipage. No enganchar
% \frame con \apptocmd: rompe la lectura de [plain] y de los titulos.

\setbeamertemplate{itemize item}{\raisebox{0.14em}{\textcolor{acento}{\rule{0.4em}{0.4em}}}}
\setbeamertemplate{itemize subitem}{\raisebox{0.16em}{\textcolor{tinta3}{\rule{0.32em}{0.32em}}}}
\setlength{\leftmargini}{1.3em}
\setlength{\parskip}{0.4em}

% Tablas: hairline del sistema y aire suficiente entre filas.
\setlength{\heavyrulewidth}{\altofilete}
\setlength{\lightrulewidth}{\hairline}
\setlength{\aboverulesep}{0.5ex}
\setlength{\belowrulesep}{0.7ex}

% ============================================================================
% Piezas
% ============================================================================

% --- Separador de parte ------------------------------------------------------
% Lamina limpia: filete, rotulo de etapa, titulo grande y una linea de encuadre.
%
% Sin [plain]: con el, estas laminas salian sin pie y sin numero, y la numeracion
% visible saltaba de 5 a 7, de 8 a 10 y asi cinco veces. En un documento que se
% cita por lamina y sobre el que se van a mandar comentarios por escrito, eso
% confunde de inmediato. El pie queda, con su numero.
\newcommand{\parte}[3]{%
  \begin{frame}
    \vfill
    \hspace{\margenlat}\begin{minipage}{0.78\paperwidth}
      \raggedright
      \reglamarca\\[1.0em]
      \rotulo{#1}\\[0.35em]
      {\fontsize{26}{30}\selectfont\textcolor{tinta}{\textbf{#2}}}\\[0.7em]
      {\small\textcolor{tinta2}{#3}}
    \end{minipage}
    \vfill
  \end{frame}
}

% --- Lamina de afirmacion ----------------------------------------------------
% Una idea sostenida por un parrafo, con mucho aire alrededor.
\newcommand{\afirmacion}[2]{%
  \vspace{0.3em}
  \begin{minipage}{0.88\textwidth}
    \raggedright
    {\fontsize{15}{19}\selectfont\textcolor{tinta}{#1}}\\[0.9em]
    {\footnotesize\textcolor{tinta2}{#2}}
  \end{minipage}}

% --- Imagen a caja fija ------------------------------------------------------
% Tres cajas y solo tres, en vez de catorce alturas escritas a mano por lamina.
% La figura se dimensiona por la caja, de modo que dos laminas seguidas no
% cambian de peso. keepaspectratio hace que mande la restriccion mas estrecha.
%
%   \imagen        ancha, ocupa el ancho util entero
%   \imagenmedia   media, para la columna mayor de una lamina a dos columnas
%   \imagenalta    alta, para las capturas en retrato, donde manda la altura
\newcommand{\imagen}[1]{%
  \begin{center}\includegraphics[width=\textwidth, height=\altoutil, keepaspectratio]{#1}\end{center}}

\newcommand{\imagenmedia}[1]{%
  \includegraphics[width=\linewidth, height=0.52\paperheight, keepaspectratio]{#1}}

\newcommand{\imagenalta}[1]{%
  \includegraphics[width=\linewidth, height=0.66\paperheight, keepaspectratio]{#1}}

% --- Cifra -------------------------------------------------------------------
% Dato y su leyenda. La hairline superior las alinea cuando van varias en fila.
\newcommand{\cifra}[2]{%
  \begin{minipage}[t]{0.22\textwidth}
    \raggedright
    \filete\\[0.55em]
    {\color{acento}\fontsize{24}{26}\selectfont\textbf{#1}}\\[0.3em]
    {\scriptsize\color{tinta2} #2}
  \end{minipage}}

% Cifra de anclaje: la misma pieza en tamano de titular, para sostener una
% lamina de \afirmacion sin inventar contenido. Lleva la cifra que la frase ya
% menciona, no un dato nuevo.
\newcommand{\cifraancla}[3][0.30\textwidth]{%
  \begin{minipage}[t]{#1}
    \raggedright
    \filete\\[0.6em]
    {\color{acento}\fontsize{22}{25}\selectfont\textbf{#2}}\\[0.45em]
    {\scriptsize\color{tinta2} #3}
  \end{minipage}}

% Rotulo con hairline superior, sin cifra: para enumerar condiciones en fila.
\newcommand{\rotuloancla}[2][0.21\textwidth]{%
  \begin{minipage}[t]{#1}
    \raggedright
    \filete\\[0.6em]
    {\small\color{tinta}\textbf{#2}}
  \end{minipage}}

% --- Chip --------------------------------------------------------------------
% Etiqueta de estado. Nodo de tikz para que la caja se ajuste al texto y quede
% alineada con la linea base.
\newcommand{\chip}[2]{%
  \tikz[baseline=(c.base)]{%
    \node[rectangle, rounded corners=2pt, inner xsep=5pt, inner ysep=2.6pt,
          fill=#1!10, text=#1, font=\scriptsize\bfseries] (c) {#2};}}

% --- Cajas de diagrama -------------------------------------------------------
% El nodo ES la caja: el texto no puede desbordarla porque la caja se
% dimensiona a partir del texto. Dibujar rectangulo y texto por separado
% permite que el texto se salga, que es el defecto que esto corrige.
% El argumento del estilo es el ancho de texto. El inner sep de 7 pt es parte de
% la rejilla de los diagramas: cambiarlo estrecha el hueco de las flechas.
\tikzset{
  caja/.style={
    rectangle, rounded corners=2pt, draw=none,
    text width=#1, align=flush left, inner sep=7pt,
    anchor=north west, minimum height=3.55cm},
  cajarot/.style={
    rectangle, rounded corners=2pt, draw=linea, line width=\hairline,
    text width=#1, align=flush left, inner sep=7pt,
    anchor=north west, minimum height=3.55cm},
}

% Contenido tipico de una caja: rotulo, titulo y descripcion.
\newcommand{\cajatexto}[3]{%
  {\scriptsize\textcolor{acento}{\textbf{\MakeUppercase{#1}}}}\\[0.35em]
  {\small\textcolor{tinta}{\textbf{#2}}}\\[0.35em]
  {\scriptsize\textcolor{tinta2}{#3}}}

% --- Nota de fuente ----------------------------------------------------------
% Siempre al pie del contenido, en gris terciario y sin justificar.
% A 7 pt y no en scriptsize: la cita de fuente competia con el cuerpo de la
% lamina cuando iba debajo de una imagen. Es el minimo del sistema, el mismo del
% pie de lamina, de modo que sigue siendo legible proyectada.
\newcommand{\fuente}[1]{\vspace{0.3em}\par{\raggedright\fontsize{7}{8.4}\selectfont\color{tinta3} #1\par}}


% xcolor define \textcolor con un argumento corto, de modo que un salto de parrafo
% dentro del texto coloreado aborta la compilacion. Aqui hay bloques de dos y tres
% parrafos en gris, asi que se declara largo. Es la misma definicion de xcolor, solo
% que \long.
\makeatletter
\long\def\@textcolor#1#2#3{\protect\leavevmode{\color#1{#2}#3}}
\makeatother

% --- Piezas propias de este guion -------------------------------------------
% Tarjeta de resultado: lo que viaja de un paso al siguiente. Es un nodo de tikz,
% de modo que la caja se dimensiona a partir del texto y no puede desbordarse.
\tikzset{
  artefacto/.style={
    rectangle, rounded corners=3pt, draw=acento, line width=0.7pt,
    fill=acentosuave, text width=#1, align=flush left, inner sep=8pt,
    anchor=north west},
  artefacto/.default=5.6cm,
}

% Lamina de transicion entre dos pasos. Siempre la misma composicion: a la
% izquierda el resultado que sale, a la derecha el paso que lo recibe, y debajo el
% parrafo que explica por que ese resultado basta.
%   #1 rotulo del resultado   #2 contenido del resultado
%   #3 rotulo del paso        #4 titulo del paso   #5 bajada del paso
%   #6 parrafo
% Sin [plain], por lo mismo que \parte: estas cuatro laminas salian sin numero.
%
% El text width del nodo de la derecha baja de 5,9 a 5,6 cm. Con 5,9 la caja se
% pasaba 2,28 pt del ancho y el log daba Overfull hbox en las cuatro laminas de
% transicion, siempre el mismo desbordamiento porque viene de este macro.
\newcommand{\transicion}[6]{%
  \begin{frame}
    \vfill
    \hspace{\margenlat}\begin{minipage}{0.88\paperwidth}
      \raggedright
      \begin{tikzpicture}
        \node[artefacto=5.5cm] (a) at (0, 0)
          {{\scriptsize\textcolor{acento2}{\textbf{\MakeUppercase{#1}}}}\\[0.35em]{\scriptsize\textcolor{tinta2}{#2}}};
        \coordinate (p) at ($(a.east) + (1.5cm, 0)$);
        \draw[-{Latex[length=2.2mm, width=1.8mm]}, acento, line width=0.9pt]
          ($(a.east) + (0.35cm, 0)$) -- ($(p) - (0.25cm, 0)$);
        \node[anchor=west, text width=5.6cm, align=flush left] at (p)
          {{\scriptsize\textcolor{tinta3}{\MakeUppercase{#3}}}\\[0.3em]{\large\textcolor{tinta}{\textbf{#4}}}\\[0.3em]{\scriptsize\textcolor{tinta2}{#5}}};
      \end{tikzpicture}

      \vspace{1.4em}
      {\small\textcolor{tinta2}{#6}}
    \end{minipage}
    \vfill
  \end{frame}}

\begin{document}

% ============================================================ 1. PORTADA
% El mapa es la red electrica del pais, lineas de transmision y subestaciones, y
% va como imagen de portada: sin leyenda, sin rotulos y sin escala. Dice de que
% trata el trabajo, que es buscar terreno donde esa red ya llega.
%
% Va en un overlay de tikz anclado al borde derecho de la pagina, no en una
% columna del flujo: asi el mapa entra entero, sin recortar el oriente del pais,
% y la columna de texto conserva su ancho. Metido en el flujo obligaba a
% estrecharla y «Prospectos Solares» se partia en dos lineas.
%
% Sin la tira de tres cifras que iba bajo el filete (21.447, 100 y 647): son
% resultados del procedimiento y no lo que la portada tiene que decir.
\begin{frame}[plain]
  \begin{tikzpicture}[remember picture, overlay]
    \node[anchor=east, inner sep=0pt]
      at ($(current page.east) - (0.55cm, 0)$)
      {\includegraphics[height=0.94\paperheight]{fig_mapa_portada.pdf}};
  \end{tikzpicture}%
  \vspace{0.2em}
  \begin{minipage}[c]{0.62\textwidth}
    \raggedright
    \isotipo{2.4em}\hspace{0.7em}%
    \raisebox{0.80em}{\begin{minipage}{6cm}
      \raggedright
      {\large\textcolor{marca}{\textbf{MÉTODOS MIXTOS}}}\\[0.1em]{\scriptsize\textcolor{tinta3}{C\,O\,N\,S\,U\,L\,T\,O\,R\,E\,S}}
    \end{minipage}}

    \vspace{1.1em}
    {\fontsize{26}{30}\selectfont\textcolor{tinta}{\textbf{Prospectos Solares}}}\\[0.45em]
    {\normalsize\textcolor{tinta2}{Identificación y caracterización de terrenos aptos\\
     para generación solar en Colombia}}

    \vspace{1.1em}
    \reglamarca

    \vspace{1.6em}
    {\small\textcolor{tinta}{Daniel Wiesner \quad$\cdot$\quad Susana Yepes\\[0.28em]
     Samuel Blanco Castellanos}}\\[0.5em]{\scriptsize\textcolor{tinta3}{SEPTIEMBRE DE 2026}}
  \end{minipage}
\end{frame}

% ============================================================ EL PROCEDIMIENTO
\section{El procedimiento}

% --- 2 ---
\begin{frame}{Punto de partida}
  \afirmacion{En Colombia el recurso solar no es el factor limitante.}
  {Lo escaso es el lote donde coinciden las cuatro condiciones a la vez. Cada una vive
  en una fuente distinta y ninguna se puede consultar a mano, uno a uno.}

  \vspace{2.0em}
  \hspace{0.01\textwidth}\begin{minipage}{0.96\textwidth}
    \raggedright
    \rotuloancla[0.22\linewidth]{Aptitud física del lote}\hfill
    \rotuloancla[0.24\linewidth]{Capacidad libre en la red}\hfill
    \rotuloancla[0.24\linewidth]{Régimen de suelo que permita el uso}\hfill
    \rotuloancla[0.22\linewidth]{Propietario identificable}
  \end{minipage}
\end{frame}

% --- 3 ---
\begin{frame}{El procedimiento, en seis pasos}{Cada paso recibe un resultado, produce otro y se lo entrega al siguiente}
  \vfill
  \begin{center}
    \includegraphics[width=\textwidth, height=0.66\paperheight, keepaspectratio]{fig_cadena.pdf}
  \end{center}
  \vfill
\end{frame}

% --- 4 ---
% Columnas en \raggedright: justificadas en 3,3 cm y sin partir palabras, la tabla
% abria rios de varios espacios («Las    grillas seleccionadas») y el log daba
% cinco Underfull hbox, dos de ellos con badness 10000. El propio estilo.tex dice
% que en columna estrecha no se justifica; esta tabla se habia quedado fuera.
%
% Sin shrink: las seis celdas de la tercera columna bajaron de dos lineas a una y
% la nota al pie de cuatro lineas a una. La tabla cabe a su tamano.
\begin{frame}{Lo que viaja entre un paso y el siguiente}{El procedimiento se puede auditar porque cada tramo deja un resultado propio}
  {\footnotesize\renewcommand{\arraystretch}{0.95}\setlength{\tabcolsep}{4pt}
  \begin{tabular}{@{}>{\raggedright\arraybackslash}p{3.35cm} >{\raggedright\arraybackslash}p{3.15cm} >{\raggedright\arraybackslash}p{6.85cm}@{}}
    \toprule
    \textcolor{tinta3}{\scriptsize PASO QUE LO PRODUCE} & \textcolor{tinta3}{\scriptsize RESULTADO QUE ENTREGA} & \textcolor{tinta3}{\scriptsize QUÉ CONTIENE}\\
    \midrule
    \textcolor{tinta2}{1.} El mapa de grillas & \textbf{El panel nacional} & 21.447 grillas de 5 por 5 km, con 43 columnas cada una.\\
    \textcolor{tinta2}{2.} Dónde se ha construido & \textbf{El perfil de referencia} & Las grillas que ya tienen planta, con sus variables medidas.\\
    \textcolor{tinta2}{3.} El modelo & \textbf{Las cien candidatas} & Las mismas 43 columnas del panel, más los tres puntajes del modelo.\\
    \textcolor{tinta2}{4.} Selección de grillas & \textbf{Las grillas seleccionadas} & Las que el cliente marca en el reporte. En esta corrida fueron diez.\\
    \textcolor{tinta2}{5.} De la grilla al lote & \textbf{La lista corta y su ficha} & Un lote por fila, con su medición y su calificación.\\
    \textcolor{tinta2}{6.} La información para comprar & \textbf{El expediente de compra} & La petición de matrículas y el certificado de tradición leído.\\
    \bottomrule
  \end{tabular}}
\end{frame}

% --- 5 ---
\begin{frame}{Reducción del universo de análisis}{Cada etapa deja fuera lo que no cumple su filtro; cuántas quedan depende de la corrida}
  \vspace{0.3em}
  \begin{center}
    \includegraphics[width=\textwidth, height=0.615\paperheight, keepaspectratio]{fig_embudo.pdf}
  \end{center}
\end{frame}

% ============================================================ PASO 2
\section{Paso 1. El mapa de grillas}

% --- 8. El paso 2 entero en una lamina --------------------------------------
% Eran dos: la de transicion, que solo servia para decir que el perfil se guarda
% hasta que haya algo con que compararlo, y la de la unidad de analisis. La
% segunda se sostiene sola y la primera se va.
%
% Se van tambien las dos que seguian, la de las familias de variables y la del
% suelo urbano: son detalle del reporte de caracterizacion, no del guion. Con
% ellas quedan sin uso fig_variables.pdf y fig_urbano.pdf.
\begin{frame}{El mapa de grillas}{Una grilla equivale a 25 km\textsuperscript{2} de territorio}
  \begin{minipage}[t]{0.60\textwidth}
    \vspace{0pt}
    \includegraphics[width=\linewidth, height=0.585\paperheight, keepaspectratio]{fig_grilla_contexto.pdf}
  \end{minipage}\hfill
  \begin{minipage}[t]{0.365\textwidth}
    \vspace{0.6em}
    \bandera\small\textcolor{tinta2}{El panel cubre \textbf{21.447 cuadrículas} de 5 por
    5 kilómetros, algo menos de la mitad del país. Grilla y celda nombran lo mismo.

    \vspace{0.9em}
    La grilla no es el proyecto. Es la unidad que hace comparable con un mismo dato todo
    lo que el panel cubre, medido a distancia, sin pedir información a nadie y sin salir a
    campo.}
  \end{minipage}
  \fuente{Celda 0008656, Montería, sobre imagen de Esri World Imagery. El panel cubre
  536.175 de los 1.168.350 km\textsuperscript{2} del país.}
\end{frame}

% ============================================================ PASO 1
\section{Paso 2. Dónde se ha construido}

% --- 8 ---
% La figura sube de 0,40 a 0,46 de la altura y el texto baja a tres lineas: las
% etiquetas del grafico estaban a 7,0 pt de interlineado, por debajo del pie de la
% propia presentacion.
\begin{frame}{Qué distingue a una grilla que ya tiene planta}{Diferencia medida frente al resto del panel, ordenada de mayor a menor}
  \vspace{-0.2em}
  \begin{center}
    \includegraphics[width=\textwidth, height=0.495\paperheight, keepaspectratio]{fig_contraste.pdf}
  \end{center}
  \vspace{-0.4em}
  \hspace{0.01\textwidth}\begin{minipage}{0.96\textwidth}
    \raggedright\footnotesize\textcolor{tinta2}{El recurso solar aparece en séptimo
    lugar: en Colombia es abundante casi en todas partes y por eso separa poco. Lo
    escaso es el acceso a la red.}
  \end{minipage}
  \fuente{A la izquierda, las grillas con al menos una planta de 10 MW o más registrada en
  el operador del mercado; a la derecha, el resto del panel. El corte es por tamaño, no por
  si ya arrancó. Cada panel lleva su escala; la privación va de 0 a 100 y más alto es peor.}
\end{frame}

% ============================================================ PASO 3
\section{Paso 3. El modelo}

% --- 8. El paso 3 en una sola lamina ----------------------------------------
% Eran cuatro: la transicion que anunciaba el paso, esta, la de los limites del
% modelo y la del radio de la subestacion. Quedan fundidas en la que explica como
% se mide el parecido, que es la unica que ensena el mecanismo.
\begin{frame}{El modelo: cómo se mide el parecido}{Dos preguntas distintas, combinadas en un solo puntaje}
  \vspace{0.2em}
  \begin{center}
    \includegraphics[width=\textwidth, height=0.54\paperheight, keepaspectratio]{fig_similitud.pdf}
  \end{center}
  \fuente{Cada componente se lleva a una escala de 0 a 1 antes de mezclarlo, porque uno es
  una distancia de Mahalanobis y el otro una euclidiana: el puntaje ordena dentro de la
  corrida y no es una nota absoluta. Usa 33 de las 43 columnas del panel.}
\end{frame}


% --- 18 ---
\transicion{Cien candidatas}
  {Identificador de celda $\cdot$ polígono $\cdot$ covariables.}
  {Paso 4}{Selección de grillas}{Por los criterios y los umbrales del cliente.}
  {El paso siguiente no necesita saber cómo se calculó el puntaje ni vuelve a abrir el
  modelo. Por eso el procedimiento es reemplazable por tramos: si mañana el modelo
  devuelve otras cien, el paso 4 corre igual.}

% ============================================================ PASO 4
\section{Paso 4. Selección de grillas}

% --- 19 ---
% Aqui habia una captura del panel de exclusiones del reporte. Es un panel de puro
% texto: puesta a 0,62 del ancho de lamina, su cuerpo quedaba por debajo de 3 pt y
% no se leia proyectada. Ademas decia «las diez plantas de 50 MW o mas» y la nota
% de esta misma lamina decia once, que es una contradiccion a la vista. Los cinco
% criterios van ahora compuestos en la lamina, legibles, y el conteo de plantas se
% enuncia de la unica forma que es cierta con los dos cortes de lectura.
\begin{frame}{Lo que impide el desarrollo}{Cuatro condiciones que se aplican antes de calificar nada}
  \vspace{0.2em}
  {\footnotesize
  \begin{tabular}{@{}>{\raggedright\arraybackslash}p{3.5cm} >{\raggedright\arraybackslash}p{10.0cm}@{}}
    \toprule
    \textbf{Figura jurídica} & Parque nacional, resguardo indígena, consejo comunitario o área protegida del RUNAP.\\
    \textbf{Altitud sobre 3.000 m} & Páramo. La Ley 1930 de 2018 prohíbe allí las actividades de alto impacto.\\
    \textbf{Cultivos de coca} & Riesgo operativo. Ninguna planta de escala utility del país está en una grilla con coca.\\
    \textbf{Conflicto armado} & Tres o más acciones bélicas desde 2022 y una densidad de tres o más por mil km\textsuperscript{2}. Ninguna planta de 50 MW o más está en un municipio así. Al corte de 10 MW, que es el que calibra los umbrales, hay cinco excepciones, tres de ellas ya en operación.\\
    \bottomrule
  \end{tabular}}

\end{frame}

% --- 20 ---
\begin{frame}{Cómo se pondera cada criterio}{La metodología es fija; los criterios y su número, no}
  \vspace{-0.3em}
  \begin{center}
    \includegraphics[width=0.94\textwidth, height=0.42\paperheight, keepaspectratio]{fig_pesos.pdf}
  \end{center}
  \vspace{0.3em}
  \hspace{0.02\textwidth}\begin{minipage}{0.94\textwidth}
    \raggedright\footnotesize\textcolor{tinta2}{El peso de cada criterio sale de una sola
    medida: cuánto separa ese criterio las grillas que ya tienen planta de las demás
    grillas al alcance de la red. Esa regla no cambia. Los siete de abajo son los que hoy
    se usan, pero el catálogo puede crecer o cambiar: al añadir un criterio, su peso se
    calcula igual y los demás se reajustan solos.}
  \end{minipage}
  \fuente{Aquí la comparación es contra las grillas al alcance de la red, no contra el panel entero.}
\end{frame}

% --- 21 ---
% Aqui habia una captura de la matriz de criterios del reporte: una tabla de siete
% filas por seis columnas capturada a 2.388 px de ancho. Colocada en la lamina, su
% cuerpo aterrizaba en unos 2 pt y no se leia ni una cifra, con lo que la unica
% lamina que ensena como se puntua un criterio no ensenaba nada. La escala se
% dibuja, con el ejemplo medido de un criterio real.
\begin{frame}{Los umbrales los pone el cliente}{Tres puntos por criterio, editables: al mover uno, las cien grillas se reclasifican a la vista}
  \vspace{0.6em}
  \begin{center}
    \includegraphics[width=\textwidth, height=0.505\paperheight, keepaspectratio]{fig_escala.pdf}
  \end{center}
  \vspace{0.6em}
  \hspace{0.02\textwidth}\begin{minipage}{0.94\textwidth}
    \raggedright\footnotesize\textcolor{tinta2}{Los tres puntos son percentiles de las 53
    celdas del país que ya tienen planta de 10 MW o más, y el cliente los reescribe. El decil
    mejor es el percentil 90 aquí, donde más es mejor.}
  \end{minipage}
\end{frame}

% --- 22 ---
% La captura del reparto tenia tres cuartas partes de blanco al pie: con
% keepaspectratio la altura mandaba y la barra salia a 3,5 cm de ancho, ilegible.
% La figura dice lo mismo, ensena ademas donde cae cada grilla en el indice y se
% lee proyectada.
\begin{frame}{Cómo se clasifica cada grilla}{El índice ordena, pero un solo criterio fuera de su límite cambia la clase}
  \vspace{0.4em}
  \begin{center}
    \includegraphics[width=0.94\textwidth, height=0.33\paperheight, keepaspectratio]{fig_distribucion.pdf}
  \end{center}

  \vspace{0.4em}
  \hspace{0.01\textwidth}\begin{minipage}{0.96\textwidth}
    \raggedright\footnotesize\textcolor{tinta2}{La clasificación no es el índice. Una
    grilla con índice alto queda \textbf{Condicionada} si un solo criterio cae por debajo
    de su límite, y las \textbf{Excluidas} no reciben calificación porque una figura
    jurídica o de riesgo las retira antes de puntuar.}
  \end{minipage}
  \fuente{Condicionada no es un rechazo: son 78 de las 100 candidatas y ocho de las diez
  grillas que pasan al paso 5. Aquí el rojo es el último escalón de la escala de la
  grilla; lo que se descarta va en gris. Cada grilla lleva escrito su motivo concreto, en
  su ficha y en la tabla del reporte.}
\end{frame}

% --- 17 ---
% El mapa era una captura de pantalla que llegaba a 3,19 cm de ancho, sin leyenda y
% con el contorno del pais casi invisible: la lamina que ensena el resultado del
% modelo y no se podia leer que significaba el color. Ahora es vectorial, con
% leyenda del indice, y la tabla de departamentos baja al pie en fila horizontal
% para dejarle el alto entero.
\begin{frame}{Dónde caen las cien candidatas ya calificadas}{Cien puntos, uno por grilla, con el color de su índice}
  \begin{minipage}[t]{0.62\textwidth}
    \vspace{0pt}
    \centering
    \includegraphics[width=\linewidth, height=0.585\paperheight, keepaspectratio]{fig_mapa.pdf}
  \end{minipage}\hfill
  \begin{minipage}[t]{0.345\textwidth}
    \vspace{0.6em}
    \bandera\footnotesize\textcolor{tinta2}{El color es el índice de la grilla, que se
    calcula en este paso y no en el anterior.

    \vspace{0.8em}
    Las cien no salen sueltas: caen en unas pocas zonas. Ahí coinciden la red cerca, el
    relieve plano y el suelo agrícola, que es lo que más pesa en el índice.}
  \end{minipage}
  \fuente{Las doce Excluidas no llevan color: una figura de riesgo las retira antes de puntuar.}
\end{frame}

% --- 15 y 16. La ficha de la grilla, en dos laminas -------------------------
% Una sola lamina no daba. El cuerpo de la ficha son 11 px sobre una pagina de
% 900: para que ese texto pase de los 7 pt del sistema hay que ensenar una
% columna de 410 px a nueve centimetros o mas, y eso no cabe junto a nada. Se
% parte en dos: aqui las dos piezas graficas, que escalan bien, y en la siguiente
% la tabla de criterios ampliada hasta que se lee.
\begin{frame}{La ficha de cada grilla}{Una página por celda, lista para imprimir o adjuntar}
  \vspace{0.1em}
  \begin{center}
    \includegraphics[width=\textwidth, height=0.60\paperheight, keepaspectratio]{fig_ficha_celda_comercial.pdf}
  \end{center}
  \fuente{Ficha de la grilla 0008656, Montería, tomada del reporte que se entrega.}
\end{frame}

% --- 16 ---
\begin{frame}{Lo que la ficha califica}{Cada criterio, con su umbral, su nota y su veredicto}
  \vspace{0.1em}
  \begin{center}
    \includegraphics[width=\textwidth, height=0.605\paperheight, keepaspectratio]{fig_ficha_criterios_grilla.pdf}
  \end{center}
  \fuente{Tres de los siete. La nota interpola entre tres anclas: límite 0, objetivo 70 y tope 100.}
\end{frame}

% --- 24 ---
\transicion{Grillas seleccionadas}
  {Identificador $\cdot$ polígono $\cdot$ datos de conexión.}
  {Paso 5}{De la grilla al lote}{Catastro, medición, criba y ficha.}
  {Aquí cambia la unidad de análisis. Hasta este punto el objeto es la celda de 25
  kilómetros cuadrados; a partir de aquí es el lote del catastro, con sus linderos, su
  código y su propietario. El número de grillas seleccionadas no altera el
  procedimiento, solo el volumen de consultas.}

% ============================================================ PASO 5
\section{Paso 5. De la grilla al lote}

% --- 25 ---
\begin{frame}{El catastro dentro de la grilla}{Se bajan todos los lotes y cada uno recibe sus propios valores}
  \begin{minipage}[t]{0.492\textwidth}
    \vspace{0pt}\centering
    \includegraphics[width=\linewidth, height=0.425\paperheight, keepaspectratio]{fig_catastro_celda.pdf}\\[0.4em]
    {\scriptsize\textcolor{tinta3}{Lo que baja del catastro}}
  \end{minipage}\hfill
  \begin{minipage}[t]{0.492\textwidth}
    \vspace{0pt}\centering
    \includegraphics[width=\linewidth, height=0.425\paperheight, keepaspectratio]{fig_catastro_clasificado.pdf}\\[0.4em]
    {\scriptsize\textcolor{tinta3}{Los mismos lotes, ya clasificados}}
  \end{minipage}

  \vspace{0.1em}
  \hspace{0.01\textwidth}\begin{minipage}{0.96\textwidth}
    \raggedright\footnotesize\textcolor{tinta2}{Se descarga el catastro completo del
    Instituto Geográfico Agustín Codazzi, sin filtrar, y cada lote se mide sobre su
    polígono.}
  \end{minipage}
  \fuente{Celda 0010961, La Unión, Sucre. Otra escala: azul idóneo, ámbar con gestión, rojo no viable.}
\end{frame}

% --- 17. De donde sale cada dato -------------------------------------------
% Faltaba. El guion contaba que el lote se mide y se clasifica, pero no de donde
% sale la informacion. Cada familia se nombra por la ENTIDAD que la publica, que
% es lo que un tercero puede verificar.
\begin{frame}{De dónde sale la información de cada lote}{Cinco familias de fuente, cada dato con la entidad que lo publica}
  \vspace{0.2em}
  \begin{center}
    \includegraphics[width=\textwidth, height=0.505\paperheight, keepaspectratio]{fig_fuentes_lote.pdf}
  \end{center}
  \fuente{Nueve entidades: Instituto de Hidrología, Meteorología y Estudios Ambientales,
  que aporta trece de las veinticinco capas, Agencia Nacional de Minería, Agencia Nacional
  de Hidrocarburos, Agencia Nacional de Tierras, Servicio Geológico Colombiano, Parques
  Nacionales Naturales, Ministerio de Ambiente, Instituto Geográfico Agustín Codazzi y
  Unidad de Planificación Rural Agropecuaria. El anexo trae qué aporta cada una.}
\end{frame}

% --- 18. Las tres clases ---------------------------------------------------
% La clase no sale de un umbral sobre el indice: sale de un catalogo de doce
% condiciones comprobadas una a una contra la fuente oficial de cada figura.
\begin{frame}{Cómo queda clasificado cada lote}{La clase no sale de un puntaje: sale de las condiciones que se comprueban}
  \vspace{0.3em}
  \begin{center}
    \includegraphics[width=\textwidth, height=0.575\paperheight, keepaspectratio]{fig_clases_lote.pdf}
  \end{center}
  \vspace{0.1em}
  \hspace{0.01\textwidth}\begin{minipage}{0.96\textwidth}
    \raggedright\footnotesize\textcolor{tinta2}{Quien lee la ficha no recibe un veredicto:
    recibe la siguiente gestión, con nombre de entidad.}
  \end{minipage}
\end{frame}

% --- 18. El filtro lo pone el cliente --------------------------------------
% Aqui iban dos laminas: la criba de 2.745 lotes a los caracterizados y la de los
% preajustes de tamano. La primera se retira; la segunda se rehace sin cifras.
%
% El cliente no necesita saber cuantos lotes cayo cada filtro, sino que los
% umbrales son suyos. La figura distingue las dos naturalezas de decision, la que
% admite grados y la que solo se cumple o no, que es lo que de verdad hay que
% entender antes de sentarse a mover controles.
\begin{frame}{El filtro lo pone el cliente}{Dos tipos de decisión: la que admite grados y la que no los admite}
  \vspace{0.3em}
  \begin{center}
    \includegraphics[width=\textwidth, height=0.52\paperheight, keepaspectratio]{fig_filtro_cliente.pdf}
  \end{center}
  \vspace{0.2em}
  \hspace{0.01\textwidth}\begin{minipage}{0.96\textwidth}
    \raggedright\footnotesize\textcolor{tinta2}{Setenta y cinco criterios en ocho grupos,
    de los que aquí van diez. Ninguno viene impuesto: el cliente fija lo que su proyecto
    necesita y la lista se rehace a la vista. Un lote sin dato no se retira: la ausencia de
    dato no es un incumplimiento.}
  \end{minipage}
\end{frame}

% --- 28 ---
% Dos laminas en lugar de una. La tabla de siete criterios por seis columnas
% numericas iba a 4,35 cm de ancho y 4,9 pt de interlineado: es el tipo de figura
% que no admite reduccion. Ahora ocupa una lamina entera.
%
% Y la tabla ya no es una captura. La que habia se tomo antes de que los umbrales
% se recalibraran a escala de lote: mostraba un indice de 77 donde el reporte que
% se entrega hoy dice 65. Se dibuja desde la salida vigente, de modo que no puede
% volver a desfasarse.
\begin{frame}{La ficha del lote}{Cada calificación queda justificada, criterio por criterio}
  \vfill
  \begin{center}
    \includegraphics[width=\textwidth, height=0.56\paperheight, keepaspectratio]{fig_ficha_criterios.pdf}
  \end{center}
  \vfill
  \hspace{0.01\textwidth}\begin{minipage}{0.96\textwidth}
    \raggedright\footnotesize\textcolor{tinta2}{El índice se reconstruye a la vista: es el
    promedio de las notas ponderado por el peso.}
  \end{minipage}
  \fuente{Vía y producción se detienen en 70: su objetivo está pegado al tope, sin recalibrar.}
\end{frame}

% --- 29 ---
% La captura de cabecera que iba aqui era anterior a la recalibracion de umbrales:
% su rotulo decia INDICE DEL LOTE 77 junto a una lamina que afirma 65. La identidad
% del lote se compone ahora con las piezas del sistema y no puede desfasarse.
\begin{frame}{Lo que la ficha dice de ese lote}{El índice califica el terreno; la clase dice qué hay que gestionar}
  \begin{minipage}[t]{0.42\textwidth}
    \vspace{0pt}
    \centering
    \includegraphics[width=\linewidth, height=0.50\paperheight, keepaspectratio]{fig_lote_ejemplo.pdf}
  \end{minipage}\hfill
  \begin{minipage}[t]{0.54\textwidth}
    \vspace{0.2em}
    \raggedright
    \rotulo{Lote del catastro}\\[0.5em]
    {\large\textcolor{tinta}{\textbf{Sabana de Torres}}}\\[0.25em]
    {\small\textcolor{tinta2}{Santander $\cdot$ 222 hectáreas}}\\[0.9em]
    \chip{aviso}{Viable con gestión}

    \vspace{0.9em}
    \cifraancla[0.42\linewidth]{65}{índice del lote sobre 100}\hfill
    \cifraancla[0.52\linewidth]{3}{condiciones que exigen trámite}
  \end{minipage}
  \fuente{Las tres: Unidad Agrícola Familiar, restitución de tierras y un humedal dentro
  del lote. La ficha nombra la entidad ante la que se adelanta cada trámite.}
\end{frame}

% ============================================================ DEL LOTE AL FOLIO
% Dos laminas antes del piloto. La cadena registral es lo que convierte un lote
% medido en un lote comprable, y hasta ahora el guion la despachaba en tres
% tarjetas al final. Va en grafico y casi sin texto: son cinco pasos y lo que hay
% que retener es el orden y quien produce cada pieza.
\section{Del lote al folio}

% --- 21 ---
\begin{frame}{Del código catastral al diagnóstico jurídico}{Cinco pasos entre el lote medido y el lote que se puede comprar}
  \vfill
  \begin{center}
    \includegraphics[width=\textwidth, height=0.515\paperheight, keepaspectratio]{fig_cadena_registro.pdf}
  \end{center}
  \vfill
  \fuente{Matrículas: Ley 1755 de 2015, sin costo, diez días hábiles. Certificado electrónico:
  23.000 pesos hasta 150 anotaciones, Resolución 2026-001726. Semáforo: Resolución 7448 de 2021.}
\end{frame}


% --- 22 ---
% El modelo NO sustituye la lectura determinista: la audita. Esa jerarquia es lo
% que hace defendible el paso, y por eso la lamina la enuncia antes que nada.
\begin{frame}{Qué añade el modelo sobre la lectura automática}{La lectura por regla es la verdad de base; el modelo la audita y amplía}
  \vspace{0.2em}
  \begin{center}
    \includegraphics[width=\textwidth, height=0.545\paperheight, keepaspectratio]{fig_diagnostico.pdf}
  \end{center}
  \vspace{0.1em}
  \hspace{0.01\textwidth}\begin{minipage}{0.96\textwidth}
    \raggedright\footnotesize\textcolor{tinta2}{Leer un folio exige seguir la cadena de
    tradición, y eso una expresión regular no lo hace. El diagnóstico se anexa a la
    ficha del lote, con la lectura por regla al lado.}
  \end{minipage}
\end{frame}

% --- Como se usa el modelo de lenguaje -------------------------------------
% La pregunta que un jefe hace primero es si se esta confiando en una IA. La
% lamina la responde con el mecanismo: no sustituye a nada, audita una lectura
% determinista que sigue siendo la verdad de base, y contesta sobre un catalogo
% cerrado de claves, de modo que no puede inventar categorias de riesgo.
\begin{frame}{Cómo se usa el modelo de lenguaje}{Cuatro tramos, y una garantía en cada uno}
  \vspace{0.2em}
  \begin{center}
    \includegraphics[width=\textwidth, height=0.585\paperheight, keepaspectratio]{fig_flujo_ia.pdf}
  \end{center}
  \fuente{El prompt y el esquema están corridos sobre el folio del piloto; falta
  publicarlos en el repositorio.}
\end{frame}

% ============================================================ EL PILOTO
% Dos laminas. La anterior termina diciendo que el terreno no es el problema y la
% red si; el piloto es la prueba de esa frase sobre un caso concreto, y la
% siguiente recoge el veredicto lote a lote.
%
% El veredicto es desfavorable y se presenta sin maquillar: un piloto que solo
% confirmara lo que se desea no probaria nada.
\section{El piloto}

% --- 31 ---
\begin{frame}{El piloto, de la grilla al lote}{Una grilla, 57 lotes medidos, un veredicto por cada uno}
  \begin{minipage}[t]{0.42\textwidth}
    \vspace{0pt}
    \centering
    \includegraphics[width=\linewidth, height=0.44\paperheight, keepaspectratio]{v5_piloto_grilla.png}
  \end{minipage}\hfill
  \begin{minipage}[t]{0.54\textwidth}
    \vspace{0.2em}
    \raggedright\footnotesize\textcolor{tinta2}{La celda \textbf{0014590}, en Melgar,
    Tolima, \textbf{no estaba entre las cien candidatas}. Se tomó del panel nacional y se
    inyectó como si hubiera pasado el filtro, para correr la cadena entera sobre un
    terreno que nadie eligió por conveniente.

    \vspace{0.9em}
    Salió \textbf{Condicionada}, con índice \textbf{39,4}. Conecta en Lanceros 115 kV a
    13,4 km, con \textbf{190,6 MW} libres en barra en 2026: la octava cifra más alta del
    país. La misma barra cae a 137,1 MW en 2027 y a 29,9 en 2028. El
    catastro cubre el \textbf{34 \%} de la grilla, y de sus 205 lotes 57 quedan
    caracterizados y 148 se apartan.}
  \end{minipage}
  \fuente{El resultado se declara sobre lo que devuelve el servicio de catastro, no sobre
  la grilla entera.}
\end{frame}

% --- 32 ---
\begin{frame}{El Poblado: el procedimiento dice que no}{Fila 21 de 57 por área catastral, índice 30 sobre 100, un criterio cumplido de siete}
  \begin{center}
    \includegraphics[width=\textwidth, height=0.44\paperheight, keepaspectratio]{fig_piloto_criterios.pdf}
  \end{center}
  \vspace{0.4em}
  \hspace{0.01\textwidth}\begin{minipage}{0.96\textwidth}
    \raggedright\footnotesize\textcolor{tinta2}{\textbf{68 \% de bosque} y 4,5 hectáreas
    aptas de 14,0, con \textbf{99,7 \% de suelo de protección} en el POT. La salida lo
    clasifica \textbf{No viable}: ninguna gestión del proyecto levanta esa condición.}
  \end{minipage}
  \fuente{\textbf{El procedimiento sirve sobre todo cuando dice que no, y dice por qué.}
  El veredicto se contrastó contra la realidad del terreno en la revisión del 27 de agosto
  y coincidió. Melgar no tiene resolución de Unidad Agrícola Familiar cargada: la ficha no
  marca exceso, pero tampoco afirma que el lote quepa.}
\end{frame}

% ============================================================ PASO 6
\section{Paso 6. La información para comprar}

% El canal entre la caja del resultado y el parrafo quedaba en menos de 2 mm y se
% leia como una colision: las tres tarjetas se reparten en 13,4 cm en lugar de
% 10,3, que dejaban un tercio de la lamina vacio a la derecha.
% El bloque de arriba era un nodo como el de las laminas de transicion, y repetia
% una composicion que el guion ya ha usado cuatro veces ademas de costar 1,7 cm de
% alto en la lamina mas cargada. Con una sola linea de texto y las descripciones
% recortadas, las tres tarjetas caben sin encoger.
\begin{frame}{Del paso 5 al paso 6, y lo que el catastro no dice}{Lo que viaja: la lista corta de lotes con su ficha}
  \hspace{0.01\textwidth}\begin{minipage}{0.96\textwidth}
    \raggedright\footnotesize\textcolor{tinta2}{Con la lista corta se sale a campo, pero
    no se firma: el propietario y la matrícula los publica la Superintendencia de
    Notariado y Registro, no el catastro.}
  \end{minipage}

  \vspace{0.1em}
  \begin{tikzpicture}[x=1cm, y=1cm]
    \node[cajarot=3.85cm, minimum height=2.45cm] (m) at (0, 0)
      {\cajatexto{Matrícula}{Automático}{El sistema arma el derecho de petición, listo
       para radicar. Diez días hábiles y silencio positivo. Gratis.}};
    \node[cajarot=3.85cm, minimum height=2.45cm] (c) at (4.72, 0)
      {\cajatexto{Certificado}{23.000 por lote}{La compra es manual, folio a folio; el
       semáforo clasifica once riesgos: falsa tradición, baldío adjudicado, restitución.}};
    \node[cajarot=3.85cm, minimum height=2.45cm] (u) at (9.44, 0)
      {\cajatexto{Unidad Agrícola Familiar}{La que decide}{Tope legal de acumulación
       sobre lo que fue baldío adjudicado. La compra que lo viola es nula: la vía es
       arriendo largo o usufructo.}};
  \end{tikzpicture}
  \fuente{Petición de matrículas: Ley 1755 de 2015. Semáforo del certificado: Resolución 7448 de 2021.}
\end{frame}

% ============================================================ CIERRE
\section{Estado}

% --- 34 ---
% «Prueba piloto sobre un lote conocido» estaba en la columna de lo pendiente el
% mismo dia en que se entrega el piloto de Melgar con sus 46 lotes: pasa a lo
% operativo. En lo pendiente queda lo que de verdad falta, que es el contraste
% contra el desempeno real de las plantas en operacion.
%
% Sin shrink: «lo que se entrega» eran tres nombres de fichero html y ahora son
% tres lineas de lenguaje corriente, que ocupan lo mismo, y las dos columnas de
% arriba perdieron una linea cada una.
\begin{frame}{Estado y siguiente paso}
  \vspace{0.1em}
  \begin{minipage}[t]{0.47\textwidth}
    \raggedright
    \rotulo[acento]{Operativo hoy}\\[0.5em]{\footnotesize\textcolor{tinta2}{Caracterización de grillas con umbrales editables.\\[0.10em]
    Descarga del catastro y medición del lote.\\[0.10em]
    Cruce con POT, valor, entorno y criba jurídica.\\[0.10em]
    Derecho de petición de matrículas y su anexo.\\[0.10em]
    Lectura automática del certificado de tradición.\\[0.10em]
    Diagnóstico jurídico del folio con modelo de lenguaje.\\[0.10em]
    Piloto completo de una grilla.}}
  \end{minipage}\hfill
  \begin{minipage}[t]{0.47\textwidth}
    \raggedright
    \rotulo[aviso]{En desarrollo}\\[0.5em]{\footnotesize\textcolor{tinta2}{Valoración por comparables de mercado.\\[0.10em]
    Automatización de la compra del certificado.\\[0.10em]
    Identificación y contacto del propietario.\\[0.10em]
    Contraste del índice contra el desempeño real de las plantas.}}
  \end{minipage}

  \vspace{-0.35em}
  \hspace{0.01\textwidth}\begin{minipage}{0.96\textwidth}
    \filete\\[0.25em]
    \begin{minipage}[t]{0.47\linewidth}
      \raggedright
      \rotulo[acento2]{Lo que se entrega}\\[0.5em]{\footnotesize\textcolor{tinta2}{Reporte de las cien candidatas.\\[0.15em]
      Visor de los 839 lotes.\\[0.15em]
      Visor del piloto, 57 lotes.}}
    \end{minipage}\hfill
    \begin{minipage}[t]{0.47\linewidth}
      \raggedright
      \rotulo[acento2]{La decisión que sigue}\\[0.5em]{\footnotesize\textcolor{tinta2}{A escala utility la barra no da: falla en 66
      de las cien. A escala distribuida sí; fijar el perfil define la lista y el costo.}}
    \end{minipage}
  \end{minipage}
\end{frame}

% ============================================================ ANEXO
\section{Anexo}

% arraystretch local: con el 1,25 del resto del documento, once filas de tabla no
% caben en la lamina y beamer la encogia un 32 %.
\begin{frame}{Fuentes y cortes}{El panel nacional y el modelo}
  {\scriptsize\renewcommand{\arraystretch}{1.14}\setlength{\tabcolsep}{4pt}
  \begin{tabular}{@{}>{\raggedright\arraybackslash}p{4.55cm} >{\raggedright\arraybackslash}p{4.95cm} >{\raggedright\arraybackslash}p{3.90cm}@{}}
    \toprule
    \textcolor{tinta3}{FUENTE} & \textcolor{tinta3}{QUÉ APORTA} & \textcolor{tinta3}{VERSIÓN O CORTE}\\
    \midrule
    Copernicus DEM GLO-30 & Pendiente, rugosidad y elevación & Release 2021, 30 m\\
    Agencia Espacial Europea & Cobertura del suelo, WorldCover & 2021 v200, 10 m\\
    Global Solar Atlas, Banco Mundial & Recurso solar en el punto & Datos 1.7, abril de 2026\\
    Unidad de Planeación Minero Energética & Capacidad libre por barra, Circular Externa 054 & 2026, preliminar, corte 5-may-2026\\
    XM, operador del mercado eléctrico & Plantas solares registradas & 253 plantas, 148 en operación\\
    Centro Nacional de Memoria Histórica & Acciones bélicas por municipio & Desde 2022\\
    OpenStreetMap & Líneas de transmisión, subestaciones y vías & Descarga 10-ago-2026, ODbL\\
    Esri World Imagery & Imagen satelital de referencia & Mosaico vivo, fecha por vista\\
    \bottomrule
  \end{tabular}}
\end{frame}

% --- Segunda y tercera lamina de fuentes. Sin siglas: los nombres de
% institucion van completos, que es como un tercero las localiza.
%
% Van en dos laminas porque en una sola no caben. El anexo nombraba cinco
% entidades y el lote se cruza contra nueve: faltaban el Instituto de
% Hidrologia, que aporta trece de las veinticinco capas, las dos agencias de
% minas e hidrocarburos y el Servicio Geologico, y esas dos primeras sostienen
% dos de las condiciones de la lamina de clases. -------------------------
\begin{frame}{Fuentes y cortes}{El lote: catastro, ordenamiento y tenencia}
  {\scriptsize\renewcommand{\arraystretch}{1.14}\setlength{\tabcolsep}{4pt}
  \begin{tabular}{@{}>{\raggedright\arraybackslash}p{4.55cm} >{\raggedright\arraybackslash}p{4.95cm} >{\raggedright\arraybackslash}p{3.90cm}@{}}
    \toprule
    \textcolor{tinta3}{FUENTE} & \textcolor{tinta3}{QUÉ APORTA} & \textcolor{tinta3}{VERSIÓN O CORTE}\\
    \midrule
    Instituto Geográfico Agustín Codazzi & Polígono del lote, código, destino y valor catastral de referencia & Corte 30-jun-2026\\
    Instituto Geográfico Agustín Codazzi & Ordenamiento territorial y capacidad de uso de las tierras & Servicio vivo; multiescalar 2024\\
    Unidad de Planificación Rural Agropecuaria & Frontera agrícola nacional & Actualización 2024\\
    Agencia Nacional de Tierras & Unidad Agrícola Familiar, resguardos y consejos comunitarios & 25-jun-2026\\
    Unidad de Restitución de Tierras & Microzonas focalizadas para restitución & 31-jul-2026\\
    Superintendencia de Notariado y Registro & Matrícula y certificado de tradición & Consulta por lote\\
    \bottomrule
  \end{tabular}}
\end{frame}

\begin{frame}{Fuentes y cortes}{El lote: el entorno}
  {\scriptsize\renewcommand{\arraystretch}{1.14}\setlength{\tabcolsep}{4pt}
  \begin{tabular}{@{}>{\raggedright\arraybackslash}p{4.55cm} >{\raggedright\arraybackslash}p{4.95cm} >{\raggedright\arraybackslash}p{3.90cm}@{}}
    \toprule
    \textcolor{tinta3}{FUENTE} & \textcolor{tinta3}{QUÉ APORTA} & \textcolor{tinta3}{VERSIÓN O CORTE}\\
    \midrule
    Instituto de Hidrología, Meteorología y Estudios Ambientales & Drenajes, antecedentes de inundación, sequía, incendios y régimen de lluvia & Trece de las veinticinco capas; servicio en vivo\\
    Ministerio de Ambiente y Desarrollo Sostenible & Páramos, humedales y cuencas ordenadas & Sistema de Información Ambiental\\
    Parques Nacionales Naturales de Colombia & Áreas protegidas del Registro Único Nacional & Servicio en vivo\\
    Agencia Nacional de Minería & Títulos y solicitudes mineras vigentes & Servicio en vivo\\
    Agencia Nacional de Hidrocarburos & Tierras en contrato o evaluación & Servicio en vivo\\
    Servicio Geológico Colombiano & Amenaza por movimientos en masa y amenaza sísmica & Escala 1:100.000\\
    \bottomrule
  \end{tabular}}
  \fuente{Cuatro capas de clima y suelo se consultan en el centroide del lote; el resto,
  contra su polígono. Cuando un servicio no responde, la ficha lo dice en vez de dar el
  criterio por cumplido.}
\end{frame}


\end{document}
